\chapter{Class Field Theory}
\begin{center}\itshape
A prime $p\equiv1\pmod 3$ can be written $x^2+27y^2$ exactly when $2$ is a cube modulo $p$.\\
No congruence on $p$ can say this.\\
So what kind of law is it?
\end{center}

Fermat knew that an odd prime is a sum of two squares exactly when $p\equiv1\pmod4$. Euler found the same kind of rule for
$x^2+2y^2$ ($p\equiv1,3\pmod8$) and $x^2+3y^2$ ($p\equiv1\pmod3$), and then got stuck on $x^2+27y^2$: the primes it represents,
$31,43,109,127,157,\dots$, follow no pattern modulo anything. Euler guessed the answer anyway, and Gauss proved it: $p=x^2+27y^2$
iff $p\equiv1\pmod3$ and $2$ is a cube modulo $p$. A condition on $p$ \emph{modulo} $p$, not modulo a fixed number.

Kronecker saw in 1853 what separates the two kinds of law. The primes $x^2+y^2$ are the primes that split completely in $\Q(i)$, an
\emph{abelian} extension of $\Q$; the primes $x^2+27y^2$ are those that split completely in $\Q(\sqrt[3]2,\omega)$, whose Galois group
$\mathfrak S_3$ is not abelian. And Kronecker announced a theorem: every abelian extension of $\Q$ lies in a cyclotomic field
$\Q(\zeta_m)$, so its splitting primes are given by congruences modulo $m$. His proof had gaps; Weber's of 1886 had others; Hilbert
gave the first complete one in 1896. Kronecker's ``liebster Jugendtraum'', the dream of his youth, was to do the same over imaginary
quadratic fields with the values of elliptic functions instead of roots of unity, and Hilbert made the general question his
twelfth problem.

What came next is \emph{class field theory}: Hilbert's conjectures (1898), Takagi's existence theorem (1920), Artin's reciprocity law
(1927), Chevalley's id\`eles (1936). It describes all abelian extensions of any number field $K$ in terms of $K$ alone: they are
classified by the generalised ideal class groups of Chapter~11, and the Frobenius element of a prime is its ideal class. This chapter
proves the Kronecker--Weber theorem in full, states the main theorems of class field theory for a general $K$ (with their proofs left
to the references, a long road of cohomology), draws their consequences for $p=x^2+ny^2$, proves cubic reciprocity and Gauss's theorem
by Gauss and Jacobi sums, and ends with one page on what happens when the Galois group is not abelian: the Langlands programme.

\begin{reminder}[from Volumes 1, 2, 5 and Chapters 4, 7, 9 and 11]
\emph{Volume~1.} $\Z[\omega]$, $\omega=e^{2\pi i/3}$, is Euclidean; its units are $\pm1,\pm\omega,\pm\omega^2$; its primes are
$\lambda=1-\omega$ (with $3=-\omega^2\lambda^2$), the rational primes $q\equiv2\pmod3$, and the factors $\pi,\bar\pi$ of the primes
$p=\pi\bar\pi\equiv1\pmod3$ (Volume~1, Theorem~10.6). $(\Z/p^k\Z)^\times$ is cyclic for $p$ odd (Volume~1, Theorem~5.5); a $2$-adic unit is a
square iff it is $\equiv1\pmod8$ (Volume~1, Corollary~12.8). \emph{Volume~2.} Galois theory; Dedekind's independence of automorphisms
(Volume~2, Lemma~10.15); characters of a finite abelian group and orthogonality (Volume~2, Section~11.4). \emph{Volume~5.} For a non-trivial
character $\chi$ modulo $p$, $\tau(\chi)=\sum_t\chi(t)e^{2\pi it/p}$ has $|\tau(\chi)|=\sqrt p$ (Volume~5, Theorem~3.5). From this book:
inertia groups and Frobenius elements (Theorem~\ref{thm:decomp}, Proposition~\ref{prop:relative}); splitting in cyclotomic fields
(Theorem~\ref{thm:cyclosplit}, Corollary~\ref{cor:cycloram}); $|d_K|>1$ for $K\ne\Q$ (Theorem~\ref{thm:dkbig}); the prime $1-\zeta$
(Proposition~\ref{prop:lambda}); unramified extensions of $p$-adic fields (Theorem~\ref{thm:unram}), Eisenstein polynomials
(Theorem~\ref{thm:eisensteinlocal}), the tower $F\subseteq T\subseteq L$ (Proposition~\ref{prop:tower}), completions
(Theorem~\ref{thm:completions}); narrow ray class groups $\mathrm{Cl}_m^+=I_m/P_m^+$ (Chapter~11) and Chebotarev's theorem
(Corollary~\ref{cor:chebcons}).
\end{reminder}

\section{The Kronecker--Weber theorem}

\begin{theorem}[Kronecker--Weber]\label{thm:kw}
Every finite abelian extension $K$ of $\Q$ is contained in a cyclotomic field $\Q(\zeta_m)$, with $m$ divisible only by primes that
ramify in $K$.
\end{theorem}

The converse is Galois theory: every subfield of $\Q(\zeta_m)$ is abelian, since
\[
\operatorname{Gal}(\Q(\zeta_m)/\Q)\cong(\Z/m\Z)^\times.
\]
We
first prove the local analogue, where there are no class groups and no units to worry about, and then glue.

\subsection*{Local Kronecker--Weber}

Throughout, $p$ is a prime, extensions of $\Q_p$ live in a fixed algebraic closure, and ``$p$-adic field'' means a finite extension of
$\Q_p$.

\begin{lemma}\label{lem:unramroot}
Let $F$ be a $p$-adic field, $u\in\mathcal O_F^\times$ and $e\ge1$ prime to $p$. Then every root of $X^e-u$ generates an unramified
extension of $F$.
\end{lemma}
\begin{proof}
Let $\mathbb F_q$ be the residue field of $F$. The reduction $X^e-\bar u$ is separable ($e\bar u\ne0$), so it splits into distinct linear
factors over some $\mathbb F_{q^f}$. The unramified extension $F_f$ of degree $f$ (Theorem~\ref{thm:unram}) has residue field $\mathbb F_{q^f}$, and by
Hensel's lemma (Theorem~\ref{thm:henselfact}) each of the $e$ simple roots of $X^e-\bar u$ lifts to a root of $X^e-u$ in $F_f$. So all $e$
roots lie in $F_f$, and each generates a subextension of the unramified extension $F_f$, which is unramified.
\end{proof}

\begin{lemma}[local cyclotomic fields]\label{lem:localcyclo}
\begin{enumerate}[label=(\alph*)]
\item $\Q_p(\zeta_{p^k})$ is totally ramified of degree $\varphi(p^k)$ over $\Q_p$, with Galois group $(\Z/p^k\Z)^\times$ acting by
$\zeta\mapsto\zeta^a$.
\item For $p\nmid n$, $\Q_p(\zeta_n)$ is unramified over $\Q_p$. For every $N$, $e(\Q_p(\zeta_N)/\Q_p)\le\varphi(p^{v_p(N)})$.
\item Let $p$ be odd, $\zeta=\zeta_p$ and $\pi=1-\zeta$. Then $-p=\pi^{p-1}w$ with $w\equiv1\pmod\pi$, and $-p$ has a $(p-1)$-th root in
$\Q_p(\zeta_p)$.
\end{enumerate}
\end{lemma}
\begin{proof}
(a) $\Phi_{p^k}(X+1)$ is Eisenstein at $p$ (Proposition~\ref{prop:lambda}), so by Theorem~\ref{thm:eisensteinlocal} $\Q_p(\zeta_{p^k})$ is
totally ramified of degree $\varphi(p^k)$. The map $\sigma\mapsto a$ with $\sigma(\zeta)=\zeta^a$ is an injective homomorphism into
$(\Z/p^k\Z)^\times$, and the orders agree.

(b) If $p\nmid n$ and $f$ is the order of $p$ modulo $n$, then $n\mid p^f-1$, so $\zeta_n\in\Q_p(\zeta_{p^f-1})$, which is unramified
(Theorem~\ref{thm:unram}). For general $N=p^aN'$ with $p\nmid N'$, $\Q_p(\zeta_N)=U(\zeta_{p^a})$ with $U=\Q_p(\zeta_{N'})$ unramified, so
$e(\Q_p(\zeta_N)/\Q_p)=e(U(\zeta_{p^a})/U)\le[U(\zeta_{p^a}):U]\le\varphi(p^a)$.

(c) By Proposition~\ref{prop:lambda}(b), $p=\prod_{a=1}^{p-1}(1-\zeta^a)=\pi^{p-1}\prod_a\frac{1-\zeta^a}{1-\zeta}$, and
$\frac{1-\zeta^a}{1-\zeta}=1+\zeta+\dots+\zeta^{a-1}\equiv a\pmod\pi$. The product is $\equiv(p-1)!\equiv-1\pmod\pi$ (Wilson), so
$w=-p/\pi^{p-1}\equiv1\pmod\pi$. By Hensel's lemma $w=w_1^{p-1}$ with $w_1\equiv1$, and $-p=(\pi w_1)^{p-1}$.
\end{proof}

\begin{lemma}[tame abelian extensions]\label{lem:tame}
Let $K/\Q_p$ be a finite abelian extension whose ramification index $e$ is prime to $p$. Then $e\mid p-1$, and $K\subseteq\Q_p(\zeta_p,\zeta_{p^f-1})$
for some $f$.
\end{lemma}
\begin{proof}
By Proposition~\ref{prop:tower}, $K$ contains the unramified $T=\Q_p(\zeta_{p^f-1})$, $f=f(K/\Q_p)$, and $K/T$ is totally ramified of degree
$e$. Note that $p$ is a uniformiser of $T$.

\emph{A radical generator.} Let $\Pi$ be a uniformiser of $K$; then $p=\Pi^eu$ with $u\in\mathcal O_K^\times$. As $K$ and $T$ have the same residue
field, $u=u_0w$ with $u_0\in\mathcal O_T^\times$ and $w\equiv1\pmod\Pi$; by Hensel's lemma $w=w_1^e$ in $K$. Then $\beta=\Pi w_1$ satisfies
$\beta^e=pc$ with $c=u_0^{-1}\in\mathcal O_T^\times$. The polynomial $X^e-pc$ is Eisenstein over $\mathcal O_T$, so it is irreducible, and
$K=T(\beta)$ by degrees.

\emph{$e\mid p-1$.} As $K/T$ is Galois, it contains the root $\zeta_e\beta$, so $\zeta_e\in K$; and $T(\zeta_e)$ is unramified over $T$
(Lemma~\ref{lem:unramroot}) and contained in the totally ramified $K$, so $\zeta_e\in T$. Let $\tau\in\operatorname{Gal}(K/T)$ with
$\tau(\beta)=\zeta_e\beta$, and $\sigma\in\operatorname{Gal}(K/\Q_p)$ restricting to the Frobenius of $T$, which acts on roots of unity of order
prime to $p$ by $\zeta\mapsto\zeta^p$ (Theorem~\ref{thm:unram}). Put $\gamma=\sigma(\beta)/\beta$. Then $\gamma^e=\sigma(c)/c\in\mathcal O_T^\times$, so
$T(\gamma)$ is unramified over $T$ and $\gamma\in T$, hence $\tau(\gamma)=\gamma$. Now
\[
\tau\sigma(\beta)=\tau(\gamma\beta)=\zeta_e\,\sigma(\beta),\qquad \sigma\tau(\beta)=\sigma(\zeta_e)\sigma(\beta)=\zeta_e^p\,\sigma(\beta).
\]
Since $\operatorname{Gal}(K/\Q_p)$ is abelian, $\zeta_e^p=\zeta_e$, so $e\mid p-1$.

\emph{Cyclotomic.} If $p=2$ then $e=1$ and $K=T$. If $p$ is odd, let $\rho=(\pi w_1')^{(p-1)/e}\in\Q_p(\zeta_p)$, where $-p=(\pi w_1')^{p-1}$ as in
Lemma~\ref{lem:localcyclo}(c); then $\rho^e=-p$. In a field containing $K$ and $\Q_p(\zeta_p)$, $(\beta/\rho)^e=-c\in\mathcal O_T^\times$, so $\beta/\rho$ lies in an
unramified extension of $T$ (Lemma~\ref{lem:unramroot}), which is unramified over $\Q_p$, hence of the form $\Q_p(\zeta_{p^{f'}-1})$
(Theorem~\ref{thm:unram}). Taking $f''$ a common multiple of $f$ and $f'$, $K=T(\beta)\subseteq\Q_p(\zeta_p,\zeta_{p^{f''}-1})$.
\end{proof}

The wild part needs Kummer theory. For a field $F$ of characteristic $0$ containing the $p$-th roots of unity and an extension $E/F$, let
$\Delta(E/F)$ be the subgroup of $F^\times/F^{\times p}$ formed by the classes of the $a\in F^\times$ that are $p$-th powers in $E$.

\begin{lemma}[Kummer]\label{lem:kummerelem}
Let $F\supseteq\mu_p$ have characteristic $0$, and $E/F$ be Galois with group $G\cong(\Z/p\Z)^r$. For $\tau\in G$ and $a=b^p$ ($b\in E$) put
$\langle\tau,a\rangle=\tau(b)/b$. This is a well-defined pairing $G\times\Delta(E/F)\to\mu_p$, multiplicative in each variable, and:
\begin{enumerate}[label=(\alph*)]
\item if $\langle\tau,a\rangle=1$ for all $a$, then $\tau=1$; so $|\Delta(E/F)|\ge p^r$;
\item if $\langle\tau,a\rangle=1$ for all $\tau$, then $a$ is trivial in $F^\times/F^{\times p}$.
\end{enumerate}
\end{lemma}
\begin{proof}
If $b'$ is another $p$-th root of $a$ then $b'=\eta b$ with $\eta\in\mu_p\subseteq F$, fixed by $\tau$; replacing $a$ by $af^p$ replaces $b$ by $bf$. So
$\langle\tau,a\rangle$ is well defined, and it lies in $\mu_p$ since its $p$-th power is $\tau(a)/a=1$. Multiplicativity in $a$ is clear; in $\tau$,
$\tau_1\tau_2(b)/b=\tau_1\bigl(\tfrac{\tau_2b}b\bigr)\tfrac{\tau_1b}b=\tfrac{\tau_2b}b\tfrac{\tau_1b}b$, as $\tfrac{\tau_2b}b\in F$.

(a) For a character $\chi:G\to\mu_p$ put $e_\chi(y)=\sum_{\tau\in G}\chi(\tau)^{-1}\tau(y)$. Then $\tau(e_\chi(y))=\chi(\tau)e_\chi(y)$, so
$e_\chi(y)^p$ is fixed by $G$ and lies in $F$: every $e_\chi(y)$ is a $p$-th root of an element of $F$. By orthogonality,
$\sum_\chi e_\chi(y)=\sum_\tau\bigl(\sum_\chi\chi(\tau)^{-1}\bigr)\tau(y)=|G|\,y$. So $E$ is spanned over $F$ by $p$-th roots $b$ of elements of $F$. If
$\tau$ fixes all of them it is the identity. Hence $G$ embeds in the group of homomorphisms $\Delta(E/F)\to\mu_p$, whose order is
$|\Delta(E/F)|$ when $\Delta(E/F)$ is finite (otherwise there is nothing to prove).

(b) Then $b$ is fixed by $G$, so $b\in F$ and $a=b^p\in F^{\times p}$.
\end{proof}

\begin{lemma}\label{lem:Up}
Let $p$ be odd, $F=\Q_p(\zeta_p)$, $\pi=1-\zeta_p$, $\mathcal O$ its valuation ring and $U_n=1+\pi^n\mathcal O$ for $n\ge1$. For $a\in(\Z/p\Z)^\times$ let
$\sigma_a\in\operatorname{Gal}(F/\Q_p)$ map $\zeta_p$ to $\zeta_p^a$.
\begin{enumerate}[label=(\alph*)]
\item $x\mapsto1+x\pi^n$ induces a group isomorphism $\mathbb F_p\cong U_n/U_{n+1}$, on which $\sigma_a$ acts as multiplication by $a^n$.
\item $U_1^p=U_{p+1}$.
\end{enumerate}
\end{lemma}
\begin{proof}
(a) $F$ is totally ramified, so its residue field is $\mathbb F_p$, and $(1+x\pi^n)(1+y\pi^n)\equiv1+(x+y)\pi^n\pmod{\pi^{n+1}}$. Moreover
$\sigma_a(\pi)=1-\zeta^a=\pi(1+\zeta+\dots+\zeta^{a-1})\equiv a\pi\pmod{\pi^2}$, and $\sigma_a(x)\equiv x\pmod\pi$ (the residue field is $\mathbb F_p$, fixed
by $\sigma_a$). So $\sigma_a(1+x\pi^n)\equiv1+a^nx\pi^n\pmod{\pi^{n+1}}$.

(b) By Lemma~\ref{lem:localcyclo}(c), $p\equiv-\pi^{p-1}\pmod{\pi^p}$, and $v_\pi(p)=p-1$. \emph{$U_1^p\subseteq U_{p+1}$.} Every $u\in U_1$ is
$(1+c\pi)u'$ with $c\in\Z$ and $u'\in U_2$. In $(1+c\pi)^p=1+pc\pi+\sum_{k=2}^{p-1}\binom pkc^k\pi^k+c^p\pi^p$ the middle terms have valuation
$\ge p-1+2=p+1$, and $pc\pi+c^p\pi^p\equiv(c^p-c)\pi^p\equiv0\pmod{\pi^{p+1}}$ since $p\mid c^p-c$. For $u'=1+y$, $v_\pi(y)\ge2$, all terms of
$(1+y)^p-1$ have valuation $\ge p+1$. \emph{$U_{p+1}\subseteq U_1^p$.} For $n\ge2$ and $x\in\mathcal O$,
$(1+x\pi^n)^p\equiv1+px\pi^n\equiv1-x\pi^{n+p-1}\pmod{\pi^{n+p}}$, because the other terms have valuations $\ge p-1+2n$ and $pn$, both $\ge n+p$.
So every $v\in U_m$ ($m\ge p+1$) is $(1+x\pi^{m-p+1})^p$ times an element of $U_{m+1}$. Iterating gives $v=\prod_{j}(1+x_j\pi^{m-p+1+j})^p\cdot
v_\infty$ with factors tending to $1$; the infinite product converges, and $v$ is the $p$-th power of an element of $U_2$.
\end{proof}

\begin{proposition}\label{prop:nop3}
For $p$ odd, no Galois extension of $\Q_p$ has Galois group $(\Z/p\Z)^3$.
\end{proposition}
\begin{proof}
Suppose $E/\Q_p$ is one. Let $F=\Q_p(\zeta_p)$; as $[F:\Q_p]=p-1$ is prime to $p$, $E\cap F=\Q_p$, so $EF/F$ is Galois with group
$G\cong(\Z/p\Z)^3$, and $EF/\Q_p$ is abelian. Let $\Delta=\Delta(EF/F)$; it has order $\ge p^3$ (Lemma~\ref{lem:kummerelem}(a)).

\emph{$\Delta$ lies in the $\omega$-eigenspace.} Let $a=b^p\in F^\times$ with $b\in EF$, and $\tilde\sigma\in\operatorname{Gal}(EF/\Q_p)$ extending $\sigma_c$. Then
$\sigma_c(a)=\tilde\sigma(b)^p$, and for $\tau\in G$, as $\tilde\sigma\tau=\tau\tilde\sigma$,
\[
\langle\tau,\sigma_c(a)\rangle=\frac{\tau\tilde\sigma(b)}{\tilde\sigma(b)}=\tilde\sigma\Bigl(\frac{\tau(b)}b\Bigr)=\langle\tau,a\rangle^c.
\]
By Lemma~\ref{lem:kummerelem}(b), $\sigma_c(a)\equiv a^c$ in $V=F^\times/F^{\times p}$, for every $c$.

\emph{That eigenspace is small.} $F^\times=\pi^\Z\times\mu_{p-1}\times U_1$ (a unit is its Teichm\"uller root of unity times an element of $U_1$),
and $\mu_{p-1}\subseteq F^{\times p}$. Let $V_U$ be the image of $U_1$ in $V$; since $U_1\cap F^{\times p}=U_1^p$ (a $p$-th root of an element of $U_1$ is
a unit, and its Teichm\"uller part is $1$), $V_U\cong U_1/U_1^p=U_1/U_{p+1}$ by Lemma~\ref{lem:Up}(b). The quotient $V/V_U$ is generated by $\pi$, on
which every $\sigma_c$ acts trivially ($\sigma_c(\pi)/\pi$ is a unit). If $v\in V$ satisfies $\sigma_c(v)=v^c$ for all $c$, its image $\bar v$ in $V/V_U$
satisfies $\bar v^{c-1}=1$; taking $c=2$ (here $p\ge3$, so $c-1$ is prime to $p$), $\bar v=1$. So the eigenspace lies in $V_U$. Filter $V_U$ by the
images of $U_1\supseteq U_2\supseteq\dots\supseteq U_{p+1}$: an element of the eigenspace lying in $U_n$ but not $U_{n+1}$ gives a non-zero vector
of $U_n/U_{n+1}$ on which $\sigma_c$ acts by $c$; by Lemma~\ref{lem:Up}(a) it also acts by $c^n$, so $c^{n-1}\equiv1$ for all $c$, i.e.\
$(p-1)\mid n-1$, i.e.\ $n\in\{1,p\}$. Hence the eigenspace has dimension at most $2$ over $\mathbb F_p$: it cannot contain $\Delta$.
\end{proof}

\begin{proposition}[$p=2$]\label{prop:Q2}
\begin{enumerate}[label=(\alph*)]
\item $\Q_2$ has exactly seven quadratic extensions $\Q_2(\sqrt d)$, $d\in\{-1,\pm2,\pm5,\pm10\}$, all contained in $\Q_2(\zeta_{24})$.
\item No cyclic extension of $\Q_2$ of degree $4$ contains $i$.
\end{enumerate}
\end{proposition}
\begin{proof}
(a) $\Q_2^\times=2^\Z\times\Z_2^\times$, and a unit is a square iff it is $\equiv1\pmod8$; so $\Q_2^\times/\Q_2^{\times2}$ has the eight classes of
$\pm1,\pm5,\pm2,\pm10$, and the seven non-trivial ones give the quadratic extensions. Now $i=\zeta_4$, $\sqrt2=\zeta_8+\zeta_8^{-1}$,
$\sqrt{-2}=i\sqrt2$, and $\sqrt{-3}=1+2\zeta_3$; since $-3\cdot5=-15\equiv1\pmod8$, $\sqrt5\in\Q_2(\sqrt{-3})$. So all seven lie in
$\Q_2(\zeta_8,\zeta_3)=\Q_2(\zeta_{24})$.

(b) Let $L\supseteq F=\Q_2(i)$ be cyclic of degree $4$ over $\Q_2$, with generator $\sigma$. Then $L=F(\sqrt\alpha)$ for some $\alpha\in F$, and
$\sigma(L)=L$ forces $F(\sqrt{\sigma\alpha})=F(\sqrt\alpha)$, so $\sigma(\alpha)=\alpha\gamma^2$ with $\gamma\in F$ (Kummer theory of quadratic
extensions). Replacing $\gamma$ by $-\gamma$ if needed, $\sigma(\sqrt\alpha)=\gamma\sqrt\alpha$, and then
$\sigma^2(\sqrt\alpha)=\sigma(\gamma)\gamma\sqrt\alpha=N_{F/\Q_2}(\gamma)\sqrt\alpha$. As $\sigma^2$ generates $\operatorname{Gal}(L/F)$, $N(\gamma)=-1$: writing
$\gamma=x+iy$, $x^2+y^2+1=0$ with $x,y\in\Q_2$. Clearing denominators gives integers, not all even, with $a^2+b^2+c^2\equiv0\pmod 8$, which is
impossible since squares are $0,1,4\pmod8$ and an odd square is $1$.
\end{proof}

\begin{theorem}[local Kronecker--Weber]\label{thm:localkw}
Every finite abelian extension of $\Q_p$ is contained in some $\Q_p(\zeta_N)$.
\end{theorem}
\begin{proof}
A finite abelian group is a product of cyclic groups of prime power order, so $K$ is the compositum of subextensions $K_i$ that are cyclic of
prime power degree, and a compositum of subfields of cyclotomic fields is in a cyclotomic field. So let $K/\Q_p$ be cyclic of degree $\ell^m$.

\emph{$\ell\ne p$.} The ramification index divides $\ell^m$, so it is prime to $p$, and Lemma~\ref{lem:tame} applies.

\emph{$\ell=p$ odd.} Let $U$ be the unramified extension of degree $p^m$ (cyclic, Theorem~\ref{thm:unram}), and $L$ the subfield of
$\Q_p(\zeta_{p^{m+1}})$ fixed by the subgroup of order $p-1$ of $(\Z/p^{m+1}\Z)^\times\cong\Z/(p-1)\times\Z/p^m$; $L/\Q_p$ is cyclic of degree $p^m$ and
totally ramified. As $U\cap L=\Q_p$, $C=UL$ has group $(\Z/p^m)^2$. Suppose $K\not\subseteq C$. Then $G=\operatorname{Gal}(KC/\Q_p)$ is abelian of exponent
dividing $p^m$ (it embeds in $\operatorname{Gal}(K/\Q_p)\times\operatorname{Gal}(C/\Q_p)$), and maps onto $\operatorname{Gal}(C/\Q_p)$ with non-trivial kernel, so
$|G|>p^{2m}$. A group of exponent $p^m$ generated by two elements has order $\le p^{2m}$, so $G$ needs at least three generators, and $G/pG$ has a
quotient $(\Z/p)^3$. Its fixed field contradicts Proposition~\ref{prop:nop3}. So $K\subseteq C\subseteq\Q_p(\zeta_{p^{p^m}-1},\zeta_{p^{m+1}})$.

\emph{$\ell=p=2$.} For $m=1$ use Proposition~\ref{prop:Q2}(a). Let $m\ge2$, $U$ the unramified extension of degree $2^m$, and
$C=U\Q_2(\zeta_{2^{m+2}})$. Since $(\Z/2^{m+2}\Z)^\times=\{\pm1\}\times\langle5\rangle\cong\Z/2\times\Z/2^m$ ($5$ has order $2^m$ modulo $2^{m+2}$, as
$5^{2^j}\equiv1+2^{j+2}\pmod{2^{j+3}}$, and $-1\notin\langle5\rangle$ since $5^k\equiv1\pmod4$), $\operatorname{Gal}(C/\Q_2)\cong\Z/2^m\times\Z/2\times\Z/2^m$. Suppose
$K\not\subseteq C$, and let $G=\operatorname{Gal}(KC/\Q_2)$, abelian of exponent dividing $2^m$, mapping onto $\operatorname{Gal}(C/\Q_2)$ with non-trivial kernel.
The fixed field of $2G$ has group $(\Z/2)^r$ and contains $2^r-1$ distinct quadratic extensions of $\Q_2$, so $r\le3$ by
Proposition~\ref{prop:Q2}(a). Write $G\cong\Z/2^{a_1}\times\Z/2^{a_2}\times\Z/2^{a_3}$ with $m\ge a_1\ge a_2\ge a_3\ge0$. The image of $2^{m-1}G$ is
$2^{m-1}\operatorname{Gal}(C/\Q_2)\cong(\Z/2)^2$, so at least two $a_i$ equal $m$; and $|G|>2^{2m+1}$ forces $a_3\ge2$. So $G$ has a quotient
$A=(\Z/4)^3$; let $E$ be the corresponding field. The fixed field of $2A$ has group $(\Z/2)^3$, so it contains all seven quadratic extensions,
among them $\Q_2(i)$, the fixed field of the kernel of some surjection $\chi:A\to\Z/2$. As $A$ is a free $\Z/4$-module, $\chi$ lifts to $\phi:A\to\Z/4$,
surjective since $\phi\bmod2=\chi\ne0$, and the fixed field of $\ker\phi$ is a cyclic quartic extension of $\Q_2$ containing $i$. This contradicts
Proposition~\ref{prop:Q2}(b). So $K\subseteq C\subseteq\Q_2(\zeta_{2^{2^m}-1},\zeta_{2^{m+2}})$.
\end{proof}

\subsection*{From local to global}

\begin{proof}[Proof of Theorem~\ref{thm:kw}]
Let $K/\Q$ be abelian. For each prime $p$ ramified in $K$ choose a prime $P$ above it. The completion $K_P$ is generated over $\Q_p$ by the roots of the
polynomial generating $K$, so it is Galois over $\Q_p$, and its automorphisms preserve $K$, so $\operatorname{Gal}(K_P/\Q_p)$ embeds in
$\operatorname{Gal}(K/\Q)$: $K_P/\Q_p$ is abelian. By Theorem~\ref{thm:localkw}, $K_P\subseteq\Q_p(\zeta_{n_p})$ for some $n_p$; let $a_p=v_p(n_p)$, which is
$\ge1$ because $K_P/\Q_p$ is ramified (Theorem~\ref{thm:completions}(c)) and $\Q_p(\zeta_n)$ is unramified for $p\nmid n$. Put $m=\prod_pp^{a_p}$ over the
ramified primes, and $L=K(\zeta_m)$, an abelian extension of $\Q$.

Let $p\mid m$ and $\mathfrak P$ a prime of $L$ above $P$. The completion $L_{\mathfrak P}=K_P(\zeta_m)$ lies in $\Q_p(\zeta_{n_p},\zeta_m)=\Q_p(\zeta_N)$ with
$v_p(N)=a_p$, so its ramification index, which is the order of the inertia group $I_p\subseteq\operatorname{Gal}(L/\Q)$ (Theorem~\ref{thm:decomp} and
Theorem~\ref{thm:completions}(c)), is at most $\varphi(p^{a_p})$ by Lemma~\ref{lem:localcyclo}(b). (In an abelian extension the inertia group depends
only on $p$.)

Let $I$ be the subgroup of $\operatorname{Gal}(L/\Q)$ generated by the $I_p$, $p\mid m$, and $M$ its fixed field. Only primes dividing $m$ ramify in $L$, and the
inertia group of $p$ in $M/\Q$ is the image of $I_p$, which is trivial; so $M/\Q$ is unramified everywhere, and $M=\Q$ by Theorem~\ref{thm:dkbig} and
Theorem~\ref{thm:ramdisc}. Hence $I=\operatorname{Gal}(L/\Q)$, and
\[
[L:\Q]\le\prod_{p\mid m}|I_p|\le\prod_{p\mid m}\varphi(p^{a_p})=\varphi(m)=[\Q(\zeta_m):\Q].
\]
Since $\Q(\zeta_m)\subseteq L$, equality holds: $L=\Q(\zeta_m)$, and $K\subseteq\Q(\zeta_m)$.
\end{proof}

Combined with the splitting law in cyclotomic fields, Kronecker--Weber is the whole of class field theory over $\Q$.

\begin{corollary}[class field theory over $\Q$]\label{cor:cftQ}
Let $K/\Q$ be abelian and $m$ with $K\subseteq\Q(\zeta_m)$; let $H_K\subseteq(\Z/m\Z)^\times$ be the subgroup corresponding to $K$. For $p\nmid m$,
$\operatorname{Frob}_p(K/\Q)$ is the image of $p\bmod m$ in $(\Z/m\Z)^\times/H_K\cong\operatorname{Gal}(K/\Q)$; in particular $p$ splits completely in $K$ iff
$p\bmod m\in H_K$. Conversely every subgroup $H$ of $(\Z/m\Z)^\times$ is $H_K$ for exactly one $K$.
\end{corollary}
\begin{proof}
$\operatorname{Frob}_p(\Q(\zeta_m)/\Q)=\sigma_p$ (Theorem~\ref{thm:cyclosplit}), and Frobenius elements restrict to Frobenius elements
(Proposition~\ref{prop:relative}). The converse is the Galois correspondence.
\end{proof}

\begin{think}
Why must the inertia groups generate everything? Because $\Q$ has no unramified extensions at all: Minkowski's bound $|d_K|>1$. Over a
field $K$ with a non-trivial class group this step fails, and the failure is measured exactly by the class group. That is why the
abelian extensions of $\Q$ are ``just'' cyclotomic, while those of $\Q(\sqrt{-5})$ include an unramified quadratic extension that no
congruence on $\Q(\sqrt{-5})$-integers can see.
\end{think}

\section{The main theorems of class field theory}

Let $K$ be a number field and $L/K$ a finite abelian extension. Let $m$ be a positive integer divisible by every rational prime below a
prime of $K$ that ramifies in $L$, and $I_m$, $P_m^+$, $\mathrm{Cl}_m^+=I_m/P_m^+$ as in Chapter~11. Every prime $\mathfrak p\in I_m$ is unramified in $L$,
so it has a Frobenius element $\operatorname{Frob}_{\mathfrak p}(L/K)\in\operatorname{Gal}(L/K)$ (one element, not just a class, since the group is abelian;
Proposition~\ref{prop:relative}). The \emph{Artin map} is the homomorphism
\[
\operatorname{Art}_{L/K}:I_m\to\operatorname{Gal}(L/K),\qquad\prod\mathfrak p^{n_{\mathfrak p}}\mapsto\prod\operatorname{Frob}_{\mathfrak p}(L/K)^{n_{\mathfrak p}}.
\]
It is surjective: the image is a subgroup $A$, and the fixed field of $A$ is an extension of $K$ in which almost all primes split completely,
which is $K$ by Chebotarev's theorem (Corollary~\ref{cor:chebcons}).

\begin{theorem}[Artin reciprocity; admitted, not proved in this series]\label{thm:artin}
For every finite abelian extension $L/K$ there is $m$ as above such that $P_m^+\subseteq\ker\operatorname{Art}_{L/K}$. Hence the Artin map induces a
surjection $\mathrm{Cl}_m^+\to\operatorname{Gal}(L/K)$: the Frobenius element of a prime depends only on its narrow ray class modulo $m$.
\end{theorem}

\begin{theorem}[existence theorem; admitted, not proved in this series]\label{thm:cftexistence}
For every $m$ and every subgroup $H$ with $P_m^+\subseteq H\subseteq I_m$, there is a unique abelian extension $L/K$, unramified at the primes not dividing $m$,
with $\ker\operatorname{Art}_{L/K}=H$. In particular $\mathrm{Cl}_m^+$ itself is the Galois group of an abelian extension, the \emph{narrow ray class field}
modulo $m$.
\end{theorem}

Over $K=\Q$ both theorems are proved: $\mathrm{Cl}_m^+(\Q)\cong(\Z/m\Z)^\times$ (an ideal prime to $m$ has a unique positive generator $\frac ab$, sent to
$ab^{-1}\bmod m$), the ray class field is $\Q(\zeta_m)$, and Corollary~\ref{cor:cftQ} says the rest. The ``positive'' in $P_m^+$ is why $\Q(\zeta_m)$,
which is complex, appears rather than its real subfield. The special case $m=1$ deserves its own name. A real place of $K$ is \emph{unramified} in $L$
if all the places of $L$ above it are real.

\begin{theorem}[Hilbert class field; admitted, not proved in this series]\label{thm:hcf}
There is a largest abelian extension $H/K$ unramified at every prime and every real place, the \emph{Hilbert class field}, and the Artin map induces
an isomorphism $\operatorname{Cl}_K\cong\operatorname{Gal}(H/K)$. In particular $[H:K]=h_K$, and if $K/\Q$ is Galois then so is $H/\Q$.
\end{theorem}

The last claim follows from maximality: for $\sigma$ an automorphism of $\overline{\Q}$, $\sigma(H)$ is the largest unramified abelian extension of
$\sigma(K)=K$, so $\sigma(H)=H$.

\begin{corollary}\label{cor:principalsplit}
Assuming Theorem~\ref{thm:hcf}: a prime $\mathfrak p$ of $K$ splits completely in $H$ if and only if it is principal.
\end{corollary}
\begin{proof}
$\mathfrak p$ is unramified in $H$, so it splits completely iff $\operatorname{Frob}_{\mathfrak p}(H/K)=1$ (Proposition~\ref{prop:relative}), iff its class is
trivial, since $\operatorname{Art}$ is an isomorphism.
\end{proof}

Figure~\ref{fig:tower} shows the smallest interesting example. In $K=\Q(\sqrt{-5})$, $h_K=2$, and $H=K(i)=\Q(i,\sqrt5)$
(Exercise~\ref{ex:hcf5}). Notice that $H$ is abelian over $\Q$, so by Kronecker--Weber it sits inside $\Q(\zeta_{20})$.

\begin{figure}[ht]
\centering
\fig{ch10-tower}
\caption{The Hilbert class field $H=\Q(i,\sqrt5)$ of $K=\Q(\sqrt{-5})$, inside $\Q(\zeta_{20})$. Edges from $\Q$ are labelled by the ramified primes: $2$ ramifies in $\Q(i)$ and
in $K$, and $5$ in $\Q(\sqrt5)$ and in $K$, each with index $2$; in the biquadratic field $H$ each has index $2$ only, so $H/K$ is
unramified.}
\label{fig:tower}
\end{figure}

\subsection*{Primes of the form \texorpdfstring{$x^2+ny^2$}{x²+ny²}}

\begin{theorem}[assuming Theorem~\ref{thm:hcf}]\label{thm:xnyy}
Let $n\ge1$ be squarefree with $n\equiv1,2\pmod4$, $K=\Q(\sqrt{-n})$ (so $\mathcal O_K=\Z[\sqrt{-n}]$), $H$ its Hilbert class field and $h=h_K$. Let
$f\in\Z[X]$ be monic of degree $h$ with a root $\alpha$ such that $H=K(\alpha)$. Then for every prime $p\nmid2n\operatorname{disc}(f)$,
\[
p=x^2+ny^2\iff\Bigl(\frac{-n}p\Bigr)=1\ \text{ and }\ f\text{ has a root modulo }p.
\]
\end{theorem}
\begin{proof}
($\Rightarrow$) $\mathfrak p=(x+y\sqrt{-n})$ has norm $p$, so $p=\mathfrak p\bar{\mathfrak p}$ splits in $K$, which means $\bigl(\frac{-n}p\bigr)=1$
(Theorem~\ref{thm:kd}). As $\mathfrak p$ is principal it splits completely in $H$ (Corollary~\ref{cor:principalsplit}): for $\mathfrak P$ above it,
$\mathcal O_H/\mathfrak P=\mathbb F_p$, and the residue of $\alpha\in\mathcal O_H$ is a root of $f$ modulo $p$.

($\Leftarrow$) Let $p=\mathfrak p\bar{\mathfrak p}$ and $\mathfrak P$ a prime of $H$ above $\mathfrak p$; it is unramified since $H/K$ is unramified and $p\nmid2n$. As
$N\mathfrak p=p$, $\phi=\operatorname{Frob}_{\mathfrak P}(H/K)$ acts as $x\mapsto x^p$ on $\mathcal O_H/\mathfrak P$. Since $[K(\alpha):K]=h$, $f$ is irreducible over $K$ and
its roots are the $h$ distinct conjugates $\tau(\alpha)$, $\tau\in\operatorname{Gal}(H/K)$, all in $\mathcal O_H$; as $p\nmid\operatorname{disc}(f)$, their residues are
the $h$ distinct roots of $\bar f$. Let $r\in\mathbb F_p$ be a root of $\bar f$, the residue of $\tau(\alpha)$ say. Then $\phi(\tau\alpha)\equiv(\tau\alpha)^p
\equiv r^p=r\equiv\tau\alpha\pmod{\mathfrak P}$, and since the residues of distinct roots are distinct, $\phi(\tau\alpha)=\tau\alpha$. But $\tau\alpha$
generates $H$ over $K$, so $\phi=1$, and $\mathfrak p$ splits completely in $H$. By Corollary~\ref{cor:principalsplit}, $\mathfrak p=(x+y\sqrt{-n})$, and
$p=N\mathfrak p=x^2+ny^2$.
\end{proof}

\begin{example}\label{exa:n14}
For $n=5$: $H=K(i)$, $f=X^2+1$, and the theorem says $p=x^2+5y^2$ iff $\bigl(\frac{-5}p\bigr)=\bigl(\frac{-1}p\bigr)=1$, iff $p\equiv1,9\pmod{20}$, as
Lagrange and Euler knew (Volume~1). For $n=14$: $h(-56)=4$, and $H=K(\alpha)$ with $\alpha=\sqrt{2\sqrt2-1}$, a root of $f=X^4+2X^2-7$ (a numerical calculation
checks with PARI that $K(\alpha)/K$ is unramified of degree $4$). So for $p\ne2,7$,
\[
p=x^2+14y^2\iff\Bigl(\frac{-14}p\Bigr)=1\text{ and }(x^2+1)^2\equiv8\pmod p\text{ is solvable}.
\]
Below $200$ these primes are $23=3^2+14$, $127=1+14\cdot3^2$, $137=9^2+14\cdot2^2$ and $151=5^2+14\cdot3^2$; for $p=23$, indeed
$(3^2+1)^2=100=4\cdot23+8$. A direct computation verifies the criterion for all primes below $10^5$.
\end{example}

\begin{think}
Why is $p=x^2+y^2$ a congruence condition but $p=x^2+14y^2$ not? For $n=1$, $h=1$ and $H=K=\Q(i)$ is abelian over $\Q$, so its split primes are
described modulo $4$ (Corollary~\ref{cor:cftQ}). For $n=14$, $H/\Q$ has Galois group the dihedral group of order $8$, which is not abelian; if the
primes $x^2+14y^2$ were the primes in some classes modulo $m$, Chebotarev's theorem would force $H\subseteq\Q(\zeta_m)$ (the argument of
Exercise~\ref{ex:nocongruence}), which is impossible.
\end{think}

For $n$ not squarefree or $n\equiv3\pmod4$, $\Z[\sqrt{-n}]$ is a non-maximal order and $H$ must be replaced by its \emph{ring class field}. For
$n=27$, the order $\Z[3\sqrt{-3}]$ of $\Q(\omega)$ has ring class field $\Q(\omega,\sqrt[3]2)$, and Theorem~\ref{thm:xnyy} becomes Gauss's theorem.
That one we can prove without any admitted theorem.

\section{Cubic reciprocity and \texorpdfstring{$x^2+27y^2$}{x²+27y²}}

Let $D=\Z[\omega]$. For a prime $\pi$ of $D$ with $N\pi\ne3$, $D/\pi D$ is a field with $N\pi$ elements (Proposition~\ref{prop:normideal}), and $N\pi\equiv1
\pmod3$: it is $p\equiv1$ or $q^2$ with $q\equiv2\pmod3$. The cube roots of unity $1,\omega,\omega^2$ stay distinct modulo $\pi$, since their
differences divide $3$.

\begin{definition}
For such $\pi$ and $\alpha\in D$, the \emph{cubic residue character} $\chi_\pi(\alpha)$ is $0$ if $\pi\mid\alpha$, and otherwise the unique cube root of
unity with
\[
\alpha^{(N\pi-1)/3}\equiv\chi_\pi(\alpha)\pmod\pi.
\]
A prime $\pi$ is \emph{primary} if $\pi\equiv2\pmod3$, that is $\pi=a+b\omega$ with $a\equiv2$, $b\equiv0\pmod3$.
\end{definition}

The power $\alpha^{(N\pi-1)/3}$ is a root of $X^3-1$ in the field $D/\pi D$, so $\chi_\pi(\alpha)$ exists and is unique. It is multiplicative in $\alpha$,
depends only on $\alpha\bmod\pi$, and (the multiplicative group of $D/\pi D$ being cyclic of order divisible by $3$) $\chi_\pi(\alpha)=1$ iff $\alpha$ is a
non-zero cube modulo $\pi$. Conjugating the congruence gives $\chi_{\bar\pi}(\bar\alpha)=\overline{\chi_\pi(\alpha)}$. For a rational prime $q\equiv2\pmod3$
and an integer $n$ prime to $q$, $n^{(q^2-1)/3}=(n^{q-1})^{(q+1)/3}\equiv1$, so $\chi_q(n)=1$.

\begin{lemma}\label{lem:primary}
A rational prime $q\equiv2\pmod3$ is primary. If $N\pi=p\equiv1\pmod3$, exactly one of the six associates of $\pi$ is primary.
\end{lemma}
\begin{proof}
The first claim is the definition. For the second, the classes $a+b\omega$ modulo $3$ with $a^2-ab+b^2\equiv1\pmod3$ are the six classes
$(a,b)\equiv(\pm1,0),(0,\pm1),\pm(1,1)$, as one checks on the nine pairs. Distinct associates $u\pi\ne u'\pi$ are incongruent modulo $3$: otherwise $3\mid(u-u')\pi$,
and $\pi$ is prime to $3$, so $3\mid u-u'$, which fails for distinct units. So the six associates of $\pi$, all of norm $p$, fall into the six classes, and
exactly one of them is $\equiv(2,0)$.
\end{proof}

From now on, let $\pi$ be primary with $N\pi=p\equiv1\pmod3$. Then $\Z/p\Z\to D/\pi D$ is an isomorphism (both have $p$ elements), and $\chi=\chi_\pi$
restricted to $\Z$ is a character of order $3$ of $(\Z/p\Z)^\times$, extended by $\chi(0)=0$; note $\chi(-1)=1$, since $\chi(-1)^3=\chi(-1)$ and
$\chi(-1)=\pm1$. Let $\zeta=e^{2\pi i/p}$, $R=\Z[\omega,\zeta]$, and
\[
g(\chi)=\sum_{t=1}^{p-1}\chi(t)\zeta^t\in R,\qquad J(\chi,\chi)=\sum_{t\bmod p}\chi(t)\chi(1-t)\in D.
\]

\begin{lemma}[Gauss and Jacobi sums]\label{lem:jacobi}
\begin{enumerate}[label=(\alph*)]
\item $g(\chi)^2=J(\chi,\chi)\,g(\overline{\chi})$ and $g(\chi)g(\overline{\chi})=p$; hence $g(\chi)^3=p\,J(\chi,\chi)$.
\item $J(\chi,\chi)\equiv-1\pmod 3$.
\item $J(\chi,\chi)=\pi$.
\end{enumerate}
\end{lemma}
\begin{proof}
(a) $g(\chi)^2=\sum_{s,t}\chi(s)\chi(t)\zeta^{s+t}=\sum_u\zeta^u\sum_{s+t=u}\chi(s)\chi(t)$. For $u\ne0$, substituting $s=us'$, $t=ut'$, the inner sum is
$\chi(u)^2J(\chi,\chi)=\overline{\chi}(u)J(\chi,\chi)$. For $u=0$ it is $\sum_s\chi(s)\chi(-s)=\chi(-1)\sum_{s\ne0}\overline{\chi}(s)=0$. So $g(\chi)^2=J(\chi,\chi)
\sum_{u\ne0}\overline{\chi}(u)\zeta^u=J(\chi,\chi)g(\overline{\chi})$. Next, $\overline{g(\chi)}=\sum\overline{\chi}(t)\zeta^{-t}=\overline{\chi}(-1)g(\overline{\chi})=
g(\overline{\chi})$, so $g(\chi)g(\overline{\chi})=|g(\chi)|^2=p$ (Volume~5, Theorem~3.5). Multiply the two identities.

(b) In any commutative ring, $(x+y)^3\equiv x^3+y^3\pmod3$. So $g(\chi)^3\equiv\sum_t\chi(t)^3\zeta^{3t}=\sum_{t\ne0}\zeta^{3t}=-1\pmod{3R}$. By (a), $pJ\equiv-1$,
and $p\equiv1\pmod3$, so $J+1\in3R$. Then $\frac{J+1}3$ is an algebraic integer in $\Q(\omega)$, hence in $D$.

(c) Let $k=\frac{p-1}3$. Since $\chi(t)\equiv t^k\pmod\pi$ for every integer $t$ (also for $t\equiv0$),
\[
J\equiv\sum_{t=0}^{p-1}t^k(1-t)^k=\sum_{j=0}^k\binom kj(-1)^j\sum_{t=0}^{p-1}t^{k+j}\pmod\pi,
\]
and $\sum_tt^i\equiv0\pmod p$ for $0<i<p-1$ (Volume~1); here $k\le k+j\le2k<p-1$. So $\pi\mid J$. By (a), $|J|^2=|g(\chi)|^4/|g(\overline{\chi})|^2=p$, so
$N(J)=p$ and $J=u\pi$ for a unit $u$. Both $J$ and $\pi$ are primary, by (b), so $u=1$ by Lemma~\ref{lem:primary}.
\end{proof}

\begin{theorem}[cubic reciprocity; Jacobi, Eisenstein]\label{thm:cubicrec}
Let $\pi_1,\pi_2$ be primary primes of $D$ of norms different from $3$ and from each other. Then $\chi_{\pi_1}(\pi_2)=\chi_{\pi_2}(\pi_1)$.
\end{theorem}
\begin{proof}
\emph{Key step.} Let $\pi$ be primary of norm $p\equiv1\pmod3$, $\chi=\chi_\pi$, $g=g(\chi)$, and let $\varpi$ be a prime of $D$ with $N=N\varpi\ne3,p$. We
work modulo $\varpi R$, a ring of characteristic $\ell$ where $N$ is a power of $\ell$, so $x\mapsto x^N$ is additive there. As $N\equiv1\pmod3$, $\chi^N=\chi$, so
\[
g^N\equiv\sum_t\chi(t)^N\zeta^{Nt}=\sum_t\chi(t)\zeta^{Nt}=\overline{\chi(N)}\,g\pmod{\varpi R},
\]
substituting $t\mapsto N^{-1}t$ modulo $p$. Since $g\,g(\overline{\chi})=p$ and $\varpi\nmid p$, $g$ is invertible modulo $\varpi R$, and $g^{N-1}\equiv\overline{\chi(N)}$. On the
other hand $g^{N-1}=(p\pi)^{(N-1)/3}\equiv\chi_\varpi(p\pi)\pmod{\varpi}$. Two cube roots of unity congruent modulo $\varpi R$ are congruent modulo $\varpi$ in $D$
(as in Lemma~\ref{lem:jacobi}(b)), hence equal. So
\begin{equation}\label{eq:keycubic}
\chi_\varpi(p)\,\chi_\varpi(\pi)=\chi_\pi(N)^{-1}.
\end{equation}

\emph{Both rational.} Then both sides of the theorem are $1$.

\emph{$\pi_1=\pi$ of norm $p$, $\pi_2=q$ rational.} Apply~\eqref{eq:keycubic} with $\varpi=q$, $N=q^2$: $\chi_q(p)=1$, so $\chi_q(\pi)=\chi_\pi(q)^{-2}=\chi_\pi(q)$.

\emph{Norms $p_1,p_2\equiv1\pmod3$.} \eqref{eq:keycubic} with $(\pi,\varpi)=(\pi_1,\pi_2)$, then $(\pi_2,\pi_1)$, gives
\[
\chi_{\pi_2}(\pi_1)=\chi_{\pi_1}(p_2)^{-1}\chi_{\pi_2}(p_1)^{-1},\qquad \chi_{\pi_1}(\pi_2)=\chi_{\pi_2}(p_1)^{-1}\chi_{\pi_1}(p_2)^{-1},
\]
and the right-hand sides are equal.
\end{proof}

\begin{theorem}[Gauss]\label{thm:gauss27}
Let $p\equiv1\pmod3$ be prime. Then $2$ is a cube modulo $p$ if and only if $p=x^2+27y^2$ with integers $x,y$.
\end{theorem}
\begin{proof}
Write $p=\pi\bar\pi$ with $\pi=a+b\omega$ primary. As $D/\pi D\cong\Z/p\Z$, $2$ is a cube modulo $p$ iff $\chi_\pi(2)=1$, iff $\chi_2(\pi)=1$ by cubic reciprocity
($2$ is a primary prime). Now $D/2D=\mathbb F_4$, whose multiplicative group has order $3$, so $\chi_2(\alpha)\equiv\alpha^{(4-1)/3}=\alpha\pmod 2$: $\chi_2(\pi)=1$
iff $\pi\equiv1\pmod2$, iff $a$ is odd and $b$ even.

($\Rightarrow$) Then $b\equiv0\pmod6$, say $b=6y$, and $4p=4(a^2-ab+b^2)=(2a-b)^2+3b^2=(2a-6y)^2+108y^2$, so $p=(a-3y)^2+27y^2$.

($\Leftarrow$) If $p=x^2+27y^2$, then $3\nmid x$; change the sign of $x$ so that $x\equiv2\pmod3$. Put $\pi=x+3y\sqrt{-3}=(x+3y)+6y\omega$ (as
$\sqrt{-3}=1+2\omega$). Its norm is $(x+3y)^2-6y(x+3y)+36y^2=x^2+27y^2=p$, so it is prime; it is primary ($x+3y\equiv2$, $6y\equiv0\pmod3$); and
$\pi\equiv x+y\pmod2$, which is odd because $p$ is. So $\chi_2(\pi)=1$, and $2$ is a cube modulo $p$.
\end{proof}

\begin{corollary}\label{cor:cuberoot2}
For a prime $p\ne2,3$: $X^3-2$ has three roots modulo $p$ iff $p=x^2+27y^2$; exactly one root iff $p\equiv2\pmod3$; and none otherwise. So the primes
$x^2+27y^2$ are exactly the primes that split completely in $\Q(\sqrt[3]2,\omega)$, and they have density $\frac16$.
\end{corollary}
\begin{proof}
If $p\equiv2\pmod3$, cubing is a bijection of $\mathbb F_p^\times$ (its kernel has order $\gcd(3,p-1)=1$), so there is exactly one root. If $p\equiv1\pmod3$,
$\mathbb F_p$ contains the cube roots of unity, so $X^3-2$ has three roots or none, and Theorem~\ref{thm:gauss27} says which. By Dedekind's theorem
(Theorem~\ref{thm:dedekindcycle}), three roots means $\operatorname{Frob}_p=1$ in $\operatorname{Gal}(\Q(\sqrt[3]2,\omega)/\Q)\cong\mathfrak S_3$, and Chebotarev gives the density
$\frac16$ (Figure~\ref{fig:cubefrob}).
\end{proof}

\begin{figure}[ht]
\centering
\fig{ch10-gauss27}
\caption{The primes $p\equiv1\pmod3$ below $1200$, in rows of $200$. Filled: $p=x^2+27y^2$ (equivalently, $2$ is a cube modulo $p$); open: the others.
About a third of them are filled, and no periodic pattern separates them (Exercise~\ref{ex:nocongruence}).}
\label{fig:gauss27}
\end{figure}

\section{Beyond abelian: the Langlands philosophy in one page}

Class field theory describes abelian extensions by objects attached to $K$ itself: ray classes, or, in Chevalley's language, characters of the id\`ele
class group $\mathbb A_K^\times/K^\times$ of Chapter~12. A Galois extension with group $\mathfrak S_3$, like $\Q(\sqrt[3]2,\omega)$, escapes this description, as
Corollary~\ref{cor:cuberoot2} shows: its split primes are not given by congruences. What could describe them?

Langlands's answer, proposed in a 1967 letter to Weil, is to replace characters (one-dimensional) by representations of higher dimension. The group
$\mathfrak S_3$ acts on the plane $x_1+x_2+x_3=0$ by permuting coordinates, a representation $\rho$ of dimension $2$, and the trace of $\rho(\operatorname{Frob}_p)$ is
$2$, $0$ or $-1$ according as $\operatorname{Frob}_p$ is the identity, a transposition or a $3$-cycle, that is, according as $X^3-2$ has three, one or no roots
modulo $p$. Now expand
\[
\eta(6z)\eta(18z)=q\prod_{n\ge1}(1-q^{6n})(1-q^{18n})=\sum_{k\ge1}a_kq^k.
\]
For every prime $p\ne2,3$, $a_p$ is exactly that trace: $a_p=2$ if $p=x^2+27y^2$, $a_p=-1$ for the other primes $p\equiv1\pmod3$, and $a_p=0$ for $p\equiv2\pmod3$. The left side is a \emph{modular form} of weight $1$ (Chapter~17), an analytic object with a symmetry group inside
$\mathrm{SL}_2(\Z)$, and the identity, a theorem of Hecke's theory, is a non-abelian reciprocity law: Frobenius elements on one side, Fourier
coefficients on the other.

The Langlands programme conjectures that every Galois representation $\rho:\operatorname{Gal}(\overline{\Q}/\Q)\to\mathrm{GL}_n(\C)$ (and more general ones, with
$\ell$-adic coefficients) corresponds in this way to an automorphic representation of $\mathrm{GL}_n(\mathbb A)$, with matching $L$-functions. For $n=1$ this
is class field theory. For $n=2$ and representations coming from elliptic curves over $\Q$ it is the modularity theorem of Wiles, Taylor, Breuil,
Conrad and Diamond, which implies Fermat's Last Theorem; Chapters~15--17 prepare its statement. In general it is open, and it is perhaps the
largest building site in mathematics.

\section{Exercises}

\begin{exercise}[$\star$]\label{ex:sqrt2plus}
Show that $\Q\bigl(\sqrt{2+\sqrt2}\bigr)$ is the real subfield of $\Q(\zeta_{16})$, that it is cyclic of degree $4$ over $\Q$, and that its conjugates are
$\pm\sqrt{2\pm\sqrt2}$.
\end{exercise}

\begin{exercise}[$\star$]\label{ex:nocubic}
Using Theorem~\ref{thm:kw} (with the condition on $m$), show that no cyclic cubic field is unramified outside $\{2,5\}$, and that exactly one cyclic cubic field
is unramified outside $7$. Find a polynomial for it.
\end{exercise}

\begin{exercise}[$\star$]\label{ex:maxabp}
Let $p$ be a prime. Show that every finite abelian extension of $\Q$ unramified outside $p$ is contained in some $\Q(\zeta_{p^k})$. What is the Galois
group of $\bigcup_k\Q(\zeta_{p^k})$ over $\Q$?
\end{exercise}

\begin{exercise}[$\star\star$]\label{ex:quadcond}
The \emph{conductor} of an abelian field $K$ is the least $m$ with $K\subseteq\Q(\zeta_m)$. Show that the conductor of a quadratic field is $|d_K|$.
\end{exercise}

\begin{exercise}[$\star\star$]\label{ex:hcf5}
Let $K=\Q(\sqrt{-5})$ and $L=K(i)$. Show that $L/K$ is unramified at every prime. Deduce from Theorem~\ref{thm:hcf} that $L$ is the Hilbert class field
of $K$, and from Theorem~\ref{thm:xnyy} that, for $p\ne2,5$, $p=x^2+5y^2$ iff $p\equiv1,9\pmod{20}$.
\end{exercise}

\begin{exercise}[$\star\star$]\label{ex:nocongruence}
Show that there is no integer $m$ such that, apart from finitely many exceptions, the primes $x^2+27y^2$ are exactly the primes in some residue classes
modulo $m$.
\end{exercise}

\begin{exercise}[$\star\star\star$]\label{ex:h23cft}
Let $K=\Q(\sqrt{-23})$, of class number $3$ (Exercise~\ref{ex:h23}), and $f=X^3-X-1$, of discriminant $-23$. Assuming Theorem~\ref{thm:hcf}, show that the
splitting field of $f$ is the Hilbert class field of $K$, and that for $p\ne23$,
\[
p=x^2+xy+6y^2\iff f\text{ has three distinct roots modulo }p.
\]
Find the smallest such prime.
\end{exercise}

\begin{beauty}
Kronecker--Weber is a statement about all abelian extensions of $\Q$, infinitely many fields, and its proof reduces it to a single computation: the
units of $\Q_p(\zeta_p)$ modulo $p$-th powers have a two-dimensional $\omega$-part. Gauss's theorem reaches beyond the abelian world, and its proof is
three lines of reciprocity once Jacobi sums are in place: the cube root of $2$ is caught by asking whether a prime of $\Z[\omega]$ is odd. And the same
coefficients that count the roots of $X^3-2$ modulo $p$ appear in an infinite product of $q$-series: arithmetic, algebra and analysis meet in one
identity, which is what the Langlands programme is about.
\end{beauty}

\section{Hints}

Try each exercise for a while before reading its hint, and read the hint before the solution.

\begin{hint}{ex:sqrt2plus}
$2\cos\frac\pi8=\zeta_{16}+\zeta_{16}^{-1}$, and $4\cos^2\frac\pi8=2+2\cos\frac\pi4$. What is $(\Z/16\Z)^\times/\{\pm1\}$?
\end{hint}
\begin{hint}{ex:nocubic}
A cyclic cubic subfield of $\Q(\zeta_m)$ corresponds to a quotient of order $3$ of $(\Z/m\Z)^\times$.
\end{hint}
\begin{hint}{ex:maxabp}
In the proof of Theorem~\ref{thm:kw}, $m$ is built from the ramified primes only.
\end{hint}
\begin{hint}{ex:quadcond}
Write $d_K$ as a product of the $p^*=(-1)^{(p-1)/2}p$ and of $-4$, $8$ or $-8$, and use Proposition~\ref{prop:gausssum}. For minimality, a prime that ramifies in
$K$ ramifies in $\Q(\zeta_m)$; and list the quadratic subfields of $\Q(\zeta_{4m'})$, $m'$ odd.
\end{hint}
\begin{hint}{ex:hcf5}
Look at the three quadratic subfields of $\Q(i,\sqrt5)$ and at the inertia groups of $2$ and $5$.
\end{hint}
\begin{hint}{ex:nocongruence}
Apply Chebotarev's theorem to $\Q(\sqrt[3]2,\omega,\zeta_m)$, and use Corollary~\ref{cor:cuberoot2}.
\end{hint}
\begin{hint}{ex:h23cft}
Show that $L/K$ is unramified at $23$ by comparing ramification indices: $f\equiv(X-10)^2(X-3)\pmod{23}$. Then follow the proof of Theorem~\ref{thm:xnyy}
with $\mathcal O_K=\Z[\frac{1+\sqrt{-23}}2]$.
\end{hint}

\section{Solutions}

\begin{solution}{ex:sqrt2plus}
\idea Recognise the number as $2\cos\frac{\pi}{8}$, then read everything off the Galois group of $\Q(\zeta_{16})$.
\writeup Let $\zeta=\zeta_{16}=e^{i\pi/8}$ and $c=\zeta+\zeta^{-1}=2\cos\frac\pi8>0$. Then $c^2=2+\zeta^2+\zeta^{-2}=2+2\cos\frac\pi4=2+\sqrt2$, so $c=\sqrt{2+\sqrt2}$.
The real subfield $\Q(\zeta)^+=\Q(c)$ is the fixed field of complex conjugation $\sigma_{-1}$, of degree $\frac{\varphi(16)}2=4$, with group
$(\Z/16\Z)^\times/\{\pm1\}$. Since $3^2=9$ and $3^4=81\equiv1$, while $9\not\equiv\pm1$, the class of $3$ has order $4$: the group is cyclic. The conjugates of $c$
are $\zeta^a+\zeta^{-a}=2\cos\frac{a\pi}8$, $a=1,3,5,7$, that is $\pm2\cos\frac\pi8$ and $\pm2\cos\frac{3\pi}8=\pm\sqrt{2-\sqrt2}$ (as $4\cos^2\frac{3\pi}8=2-\sqrt2$).
\end{solution}

\begin{solution}{ex:nocubic}
\idea By Kronecker--Weber the field lives in $\Q(\zeta_m)$ with $m$ built from the allowed primes, so everything is about quotients of $(\Z/m\Z)^\times$.
\writeup If $K$ is cyclic cubic and unramified outside $\{2,5\}$, then $K\subseteq\Q(\zeta_m)$ with $m=2^a5^b$, so $\operatorname{Gal}(K/\Q)\cong\Z/3$ is a quotient of
$(\Z/2^a)^\times\times(\Z/5^b)^\times$, whose order $2^{a-1}\cdot4\cdot5^{b-1}$ (or less) is prime to $3$: impossible. If $K$ is unramified outside $7$, then
$K\subseteq\Q(\zeta_{7^b})$ and $(\Z/7^b\Z)^\times$ is cyclic of order $6\cdot7^{b-1}$, with a unique subgroup of index $3$. So $K$ is unique: it is the cubic
subfield of $\Q(\zeta_7)$, generated by $\zeta_7+\zeta_7^{-1}=2\cos\frac{2\pi}7$, a root of $X^3+X^2-2X-1$. (It is cyclic and ramified only at $7$, so it exists.)
\end{solution}

\begin{solution}{ex:maxabp}
\idea The cyclotomic field produced by the proof only involves the ramified primes.
\writeup By Theorem~\ref{thm:kw}, $K\subseteq\Q(\zeta_m)$ where $m$ is divisible only by ramified primes, here only by $p$: $m=p^k$ (or $m=1$, $K=\Q$). The Galois
group of $\bigcup_k\Q(\zeta_{p^k})$ is the inverse limit of the groups $(\Z/p^k\Z)^\times$ under reduction, that is $\Z_p^\times$: an automorphism is given by
compatible exponents $a_k\in(\Z/p^k\Z)^\times$ with $\sigma(\zeta_{p^k})=\zeta_{p^k}^{a_k}$, which is a $p$-adic unit.
\end{solution}

\begin{solution}{ex:quadcond}
\idea Build $\sqrt{d_K}$ out of Gauss sums; conversely, ramification forces the primes of $d_K$ into $m$, and the $2$-part needs a closer look.
\writeup Let $d_K=\prod_ip_i^*\cdot\varepsilon$ with distinct odd primes $p_i$ and $\varepsilon\in\{1,-4,8,-8\}$ (this is the standard form of a fundamental
discriminant: if $d\equiv1\pmod4$, $d_K=d=\prod p_i^*$ since both sides are $\equiv1\pmod4$ and have the same absolute value; if $d\equiv2,3\pmod4$, $d_K=4d$ and
the $2$-part is $-4$, $8$ or $-8$). Now $\sqrt{p^*}\in\Q(\zeta_p)$ (Proposition~\ref{prop:gausssum}), $\sqrt{-4}=2i\in\Q(\zeta_4)$, $\sqrt{\pm8}=2\sqrt{\pm2}\in\Q(\zeta_8)$.
So $\sqrt{d_K}\in\Q(\zeta_{|d_K|})$, and $K=\Q(\sqrt{d_K})$.

Conversely let $K\subseteq\Q(\zeta_m)$; we may assume $m\not\equiv2\pmod4$, since $\Q(\zeta_{2m'})=\Q(\zeta_{m'})$ for $m'$ odd. Every odd $p\mid d_K$ ramifies in $K$,
hence in $\Q(\zeta_m)$, so $p\mid m$ (Corollary~\ref{cor:cycloram}). If $d_K$ is odd we are done: $|d_K|$ divides $m$. If $d_K$ is even, $2$ ramifies, so $4\mid m$.
Suppose $8\mid d_K$ but $8\nmid m$: $m=4m'$, $m'$ odd. The quadratic subfields of $\Q(\zeta_{4m'})$ correspond to the characters of order $2$ of
$(\Z/4)^\times\times(\Z/m')^\times$, which are products of $\chi_{-4}$ and of quadratic characters modulo the $p\mid m'$; the corresponding fields are the
$\Q(\sqrt{\pm\prod p^*})$ (by the first part, these fields are in $\Q(\zeta_{4m'})$, and there are as many of them as characters of order 2). None contains $\sqrt{\pm2}\cdot$(odd):
the discriminant of $\Q(\sqrt{\pm\prod p^*})$ is $\prod p^*$ or $-4\prod p^*$, never divisible by $8$. So $8\mid m$. In all cases $|d_K|$ divides $m$, so $m\ge|d_K|$.
\end{solution}

\begin{solution}{ex:hcf5}
\idea In a biquadratic field each ramified prime has inertia group of order at most $2$, because some quadratic subfield is unramified at it.
\writeup $L=\Q(i,\sqrt5)$ has group $(\Z/2)^2$ and quadratic subfields $\Q(i)$, $\Q(\sqrt5)$, $K=\Q(\sqrt{-5})$. Only $2$ and $5$ ramify in these, hence in $L$. At $2$:
$\Q(\sqrt5)$ is unramified at $2$ ($5\equiv1\pmod4$), so the inertia group $I_2\subseteq\operatorname{Gal}(L/\Q)$ fixes $\Q(\sqrt5)$ and has order $\le2$; as $2$
ramifies in $K$ with index $2$, $e_2(L/\Q)=2=e_2(K/\Q)$, and $L/K$ is unramified above $2$. At $5$: $\Q(i)$ is unramified at $5$, and the same argument applies.
No other prime ramifies. $K$ has no real places. So $L/K$ is an unramified abelian extension of degree $2=h_K$ (Chapter~5), hence $L\subseteq H$ and $L=H$
by degrees. Theorem~\ref{thm:xnyy} with $n=5$, $f=X^2+1$, $h=2$ gives: for $p\ne2,5$, $p=x^2+5y^2$ iff $\bigl(\frac{-5}p\bigr)=1$ and $\bigl(\frac{-1}p\bigr)=1$, iff
$p\equiv1\pmod4$ and $\bigl(\frac5p\bigr)=1$, iff $p\equiv1\pmod4$ and $p\equiv\pm1\pmod5$ (quadratic reciprocity), iff $p\equiv1,9\pmod{20}$.
\end{solution}

\begin{solution}{ex:nocongruence}
\idea A congruence description would make the Frobenius in $\mathfrak S_3$ a function of $p\bmod m$; Chebotarev produces primes contradicting that.
\writeup Let $L=\Q(\sqrt[3]2,\omega)$ and suppose that, outside a finite set, $p=x^2+27y^2$ iff $p\bmod m\in A$ for some set $A$ of classes. By
Corollary~\ref{cor:cuberoot2}, for $p\nmid6m$ outside that set, $\operatorname{Frob}_p(L/\Q)=1$ iff $p\bmod m\in A$. Let $N=L(\zeta_m)$ and
$G=\operatorname{Gal}(N/\Q)$; for $p\nmid6m$, $\operatorname{Frob}_p(N/\Q)$ restricts to $\operatorname{Frob}_p(L/\Q)$ and to $p\bmod m$ (Proposition~\ref{prop:relative},
Theorem~\ref{thm:cyclosplit}). Take $\sigma\in\operatorname{Gal}(N/\Q(\zeta_m))$ whose restriction to $L$ is not the identity; it exists unless $L\subseteq\Q(\zeta_m)$, which is
impossible since $L$ is not abelian. By Chebotarev (Corollary~\ref{cor:chebcons}) there are infinitely many $p$ with $\operatorname{Frob}_p(N/\Q)=\sigma$ (as a class,
and $\sigma$'s conjugates also restrict to the identity on $\Q(\zeta_m)$ and non-trivially on $L$), and infinitely many with $\operatorname{Frob}_p(N/\Q)=1$. The first
are $\equiv1\pmod m$ with $\operatorname{Frob}_p(L/\Q)\ne1$, the second are $\equiv1\pmod m$ with $\operatorname{Frob}_p(L/\Q)=1$. So the class $1\bmod m$ is both in $A$ and
not in $A$: a contradiction.
\end{solution}

\begin{solution}{ex:h23cft}
\idea The splitting field $L$ of $f$ is an $\mathfrak S_3$-field containing $K$; it is unramified over $K$ because $23$ already has index $2$ in $K$. Then the proof of
Theorem~\ref{thm:xnyy} runs with the form $x^2+xy+6y^2=N(x+y\frac{1+\sqrt{-23}}2)$.
\writeup $f$ is irreducible (no root modulo $2$, and degree $3$), and its discriminant $-4(-1)^3-27=-23$ is not a square, so $\operatorname{Gal}(L/\Q)\cong\mathfrak S_3$ and the
quadratic subfield of $L$ is $\Q(\sqrt{-23})=K$ (Volume~2). Since $\operatorname{disc}(f)=-23$ is squarefree, $\Z[\alpha]$ is the ring of integers of $\Q(\alpha)$ (Proposition~\ref{prop:index}),
so only $23$ ramifies in $\Q(\alpha)$, hence in its conjugates and in $L$. Modulo $23$, $f\equiv(X-10)^2(X-3)$ (check: $10+10+3=23$, $f(10)=989=23\cdot43$), so
$23=\mathfrak q_1^2\mathfrak q_2$ in $\Q(\alpha)$ (Theorem~\ref{thm:kd}). Ramification indices multiply in the tower $\Q\subseteq\Q(\alpha)\subseteq L$, so
$e_{23}(L/\Q)=2e'$ with $e'=e(L/\Q(\alpha))\in\{1,2\}$; and $e_{23}(L/\Q)$ divides $[L:\Q]=6$ (as $efg=6$). So $e_{23}(L/\Q)=2=e_{23}(K/\Q)$.
Hence $L/K$ is unramified at every prime, and at infinity since $K$ is imaginary. It is abelian (cyclic of order $3$), so $L\subseteq H$, and $[H:K]=3$ gives $L=H$.

Now let $p\ne23$ ($p=2$ and $p=3$ are not of the form, and $f$ has no root modulo $2$ or $3$; take $p\ge5$). If $p=x^2+xy+6y^2=N(x+y\theta)$, $\theta=\frac{1+\sqrt{-23}}2$,
then $\mathfrak p=(x+y\theta)$ is a principal prime of degree $1$, so it splits completely in $H$, and so does $p$ (it splits in $K$, and both primes above $p$ are principal,
being conjugate). Conversely if $p$ splits completely in $H$, then a prime $\mathfrak p$ of $K$ above $p$ has degree $1$ and splits completely in $H$, so it is principal,
$\mathfrak p=(x+y\theta)$ and $p=x^2+xy+6y^2$. Finally, $p$ splits completely in $H=L$ iff it splits completely in $\Q(\alpha)$ (Exercise~\ref{ex:splitdensity}: $L$ is the
Galois closure), iff $f$ has three distinct roots modulo $p$ (Theorem~\ref{thm:kd}, $p\nmid23$). The smallest such prime is $59=5^2+5\cdot2+6\cdot2^2$, and indeed
$X^3-X-1\equiv(X-4)(X-13)(X-42)\pmod{59}$.
\end{solution}

\section*{Sources and further reading}

\begin{itemize}
\item D.~A.~Cox, \emph{Primes of the Form $x^2+ny^2$}, 3rd ed., AMS Chelsea, 2022. The whole story of this chapter, from Fermat to the ring class
field, with $x^2+14y^2$ and $x^2+27y^2$ worked out.
\item K.~Ireland and M.~Rosen, \emph{A Classical Introduction to Modern Number Theory}, 2nd ed., Springer GTM~84, 1990, Chapter~9. Cubic reciprocity
by Jacobi sums, the proof followed here.
\item L.~C.~Washington, \emph{Introduction to Cyclotomic Fields}, 2nd ed., Springer GTM~83, 1997, Chapter~14. The proof of Kronecker--Weber through
the local theorem.
\item F.~Lemmermeyer, \emph{Reciprocity Laws: From Euler to Eisenstein}, Springer, 2000. History and many proofs of cubic reciprocity.
\item J.~Neukirch, \emph{Algebraic Number Theory}, Springer, 1999, Chapters~IV and~VI; J.~S.~Milne, \emph{Class Field Theory}, course notes,
\texttt{jmilne.org}. Complete proofs of Artin reciprocity and of the existence theorem.
\item J.~W.~S.~Cassels and A.~Fr\"ohlich (eds.), \emph{Algebraic Number Theory}, 2nd ed., London Math.\ Soc., 2010. Serre's and Tate's chapters on
local and global class field theory, and the famous exercises.
\item D.~Hilbert, \emph{Die Theorie der algebraischen Zahlk\"orper} (Zahlbericht), Jahresber.\ DMV 4 (1897); English translation, Springer, 1998.
\item S.~Gelbart, ``An elementary introduction to the Langlands program'', \emph{Bull.\ AMS} 10 (1984), 177--219.
\item \emph{History.} Euler's conjecture on $x^2+27y^2$, Gauss's proof, Kronecker's 1853 announcement and his
``Jugendtraum'' letter to Dedekind (1880), and the proofs of Weber (1886) and Hilbert (1896): Cox's book above, Chapters~1 and~2.
\end{itemize}
